\section{Modelos de Propagación de Malware: De Enfoques Compartimentales a Simulaciones Espaciales Basadas en Agentes y BDI}
\label{sec:estado_arte}
El estudio formal de la propagación de software malicioso ha evolucionado de forma paralela al desarrollo de las redes de comunicación y a la complejidad intrínseca de las amenazas. Para diseñar contenciones efectivas contra el ransomware moderno, resulta crucial entender cómo las metodologías de modelado han transitado desde formulaciones matemáticas abstractas hacia simulaciones espaciales, individuales y cognitivas de alta fidelidad.

\subsection{Los Modelos Epidemiológicos Clásicos de Compartimentos}
El origen del modelado de malware informático se remonta a los trabajos pioneros de Kephart y White a principios de la década de 1990 \cite{kephart1991directed, kephart1993measuring}. Inspirándose en la epidemiología biológica, adaptaron los modelos compartimentales clásicos descritos por Kermack y McKendrick en 1927. Estos modelos dividen a la población de nodos en compartimentos según su estado de salud y susceptibilidad: Susceptibles ($S$), Infectados ($I$) y Recuperados o Inmunes ($R$).
En sus formulaciones iniciales, se asumía una topología de mezcla homogénea, en la cual cualquier nodo tiene la misma probabilidad de entrar en contacto con cualquier otro dispositivo de la red. Bajo este enfoque, la dinámica de la infección se rige por sistemas de ecuaciones diferenciales ordinarias que determinan la velocidad a la que los nodos transitan de un estado a otro en función de tasas constantes de transmisión ($\beta$) y de curación o parcheo ($\mu$):

$$\frac{dS}{dt} = -\beta S I$$
$$\frac{dI}{dt} = \beta S I - \mu I$$
$$\frac{dR}{dt} = \mu I$$

Si bien estas herramientas matemáticas iniciales permitieron comprender nociones fundamentales como el umbral epidémico (el valor por encima del cual una amenaza de malware se convierte en una infección descontrolada), presentaban limitaciones severas al aplicarse a redes reales. Los modelos compartimentales homogéneos ignoran la estructura de interconexión física, asumiendo que un equipo de escritorio está en contacto directo y equivalente con un servidor de dominio, un cortafuegos perimetral o una fuente externa en Internet, lo cual físicamente es incorrecto.

\subsection{Evolución hacia Modelos de Redes Heterogéneas y de Nivel de Nodo}
A medida que la teoría de redes complejas maduró a finales del siglo XX y principios del XXI, los investigadores reemplazaron la suposición de mezcla homogénea por topologías de red explícitas y heterogéneas, tales como redes de mundo pequeño (Small-World) y redes libres de escala (Scale-Free). Estas estructuras representan de mejor manera las redes informáticas reales, caracterizadas por una ley de potencias en la distribución de grados de conectividad, donde coexisten numerosos nodos con pocas conexiones (estaciones de trabajo individuales) y unos pocos nodos centralizados hiperconectados o hubs (switches de distribución, routers y servidores centrales).

En este ámbito, Pastor-Satorras y Vespignani \cite{pastorsatorras2001epidemic} demostraron de manera analítica que en redes complejas libres de escala, el umbral epidémico desaparece o se reduce drásticamente, lo que significa que cualquier infección persistente, por pequeña que sea su tasa de contagio, es capaz de propagarse por toda la red. Este comportamiento se debe a la existencia de los hubs hiperconectados, los cuales actúan como amplificadores exponenciales de la infección lateral, justificando por qué el switch de distribución central en nuestro modelo representa un punto crítico de vulnerabilidad y control de seguridad.

Por su parte, Chernikova et al. \cite{chernikova2023modeling} sostienen que los modelos deterministas tradicionales SIR no logran capturar la latencia del malware moderno en redes estructuradas complejas. Proponen, en su lugar, la introducción de estados más detallados, como el estado de Infección Latente o Silenciosa ($ID$), dando origen a modelos como el SIIDR (Susceptible-Infectado-Latente-Recuperado) para simular amenazas que permanecen durmientes o que realizan escaneos sigilosos antes de detonar.
Adicionalmente, Yang, Liu y Zhang \cite{yang2019modeling} desarrollaron un método dinámico a nivel de nodo para modelar la propagación de ransomware. A diferencia de las ecuaciones epidemiológicas globales que promedian el estado de la población, este enfoque calcula la probabilidad individual de que cada nodo se infecte en función de la condición exacta de sus vecinos directos en la matriz de adyacencia de la red, demostrando que la topología local dicta de forma crítica la velocidad inicial del brote.

\subsection{Agentes Racionales y la Arquitectura BDI en Ciberseguridad}
A pesar de las mejoras introducidas por las redes heterogéneas, los modelos matemáticos y estocásticos a nivel de nodo continuaban presentando dificultades para representar la toma de decisiones individuales de los nodos, el comportamiento dinámico y las respuestas descentralizadas de defensa activa frente a incidentes. Para superar esto, la ciencia de la simulación se orientó hacia los Modelos de Agentes Inteligentes y, en especial, a la arquitectura cognitiva BDI (Belief-Desire-Intention) propuesta por Rao y Georgeff \cite{rao1995bdi}.

La arquitectura BDI representa de manera explícita los procesos cognitivos humanos o de agentes de software racionales a través de tres pilares:
\begin{itemize}
    \item \textbf{Beliefs (Creencias):} La información y el conocimiento local que el agente posee sobre su estado interno y sobre el entorno (por ejemplo, vulnerabilidades conocidas o amenazas detectadas en su segmento).
    \item \textbf{Desires (Deseos):} Los objetivos que el agente desea alcanzar (por ejemplo, sobrevivir a la infección, inmunizarse o propagar el ataque si es malicioso).
    \item \textbf{Intentions (Intenciones):} Los planes de acción concretos a los que el agente se ha comprometido de manera inmediata (por ejemplo, aplicar un parche o desconectar un enlace).
\end{itemize}

El uso de la racionalidad BDI descentralizada en ciberseguridad permite modelar cómo cada terminal actúa de forma autónoma e inteligente al percibir anomalías a su alrededor, sin requerir una coordinación global. Esto hace al sistema robusto frente a comportamientos de cascada de red e ideal para simular el comportamiento de un parque informático heterogéneo donde cada host se defiende en función de su riesgo percibido.

\subsection{Simulación Basada en Agentes (ABM) para la Ciberseguridad Activa}
En el modelado basado en agentes (ABM), cada dispositivo de red se modela como una de estas entidades autónomas dotada de atributos y reglas de comportamiento independientes. Las interacciones locales entre estos agentes dan origen a patrones macroscópicos emergentes que son difíciles de predecir mediante métodos analíticos puros.
Padilla et al. \cite{padilla2025serdux} aplicaron exitosamente esta metodología combinada en el modelo SERDUX-MARCIM para simular ciberataques dinámicos en infraestructuras críticas marítimas, demostrando que la integración de simulaciones multiagente con modelos epidemiológicos locales resulta ideal para evaluar el impacto sistémico de amenazas persistentes y ransomware.
Asimismo, Kato et al. \cite{kato2025agent} desarrollaron simulaciones basadas en agentes para evaluar los riesgos de propagación de ransomware en el sistema financiero japonés. Su modelo representó de forma explícita el comportamiento de los atacantes que buscan de forma lateral infectar servidores de bases de datos mediante escalamiento de privilegios, y cómo las defensas dinámicas en los hosts pueden ralentizar la infección. En el ámbito de las redes LAN locales, Cui et al. \cite{cui2021agent} exploraron cómo la interacción directa entre las estaciones de trabajo y las estrategias de defensa activa (como la detección temprana y el bloqueo automático de puertos) mitiga las fases de movimiento lateral de los ataques antes de comprometer los directorios activos compartidos.

\subsection{Modelos Espaciales Explícitos e Integración de Sistemas GIS}
Una de las fronteras más recientes en la simulación multiagente aplicada a la ciberseguridad es la incorporación del espacio geográfico o físico en el modelo. Las redes de datos universitarias o corporativas no existen en un vacío abstracto; están distribuidas en salas, edificios y campus. El trazado físico del cableado y la proximidad de los equipos imponen restricciones topológicas que condicionan el comportamiento de propagación y facilitan o dificultan el aislamiento de áreas infectadas.

La plataforma GAMA Platform, desarrollada por Taillandier et al. \cite{taillandier2019gama, amouroux2009gama}, se ha posicionado como una herramienta de referencia para este propósito, gracias a su capacidad nativa para consumir datos espaciales y georreferenciados en formato vectorial (Shapefiles creados en herramientas GIS como QGIS). Esto permite a los modelos definir la movilidad de agentes físicos y redes de comunicaciones directamente sobre planos reales de infraestructura. 
Gómez, Ramos y Santos \cite{gomez2024paradigm} destacan que la integración de variables GIS con modelos epidemiológicos basados en agentes permite simular brotes infecciosos (tanto biológicos como computacionales) en escenarios espaciales explícitos de etapas tempranas, brindando a los investigadores un mapa dinámico del avance del virus en el espacio geográfico de la organización.

\subsection{Estrategias de Mitigación y Contención Reactiva}
El fin último del modelado de propagación de malware es evaluar la efectividad de las contenciones. Al-Rimy, Saad y Portokalidis \cite{al2020ransomware} clasifican los mecanismos de defensa frente al ransomware en preventivos, detectivos y reactivos. En redes locales académicas, las medidas preventivas (como el parcheo periódico del puerto SMB TCP 445 sugerido por Microsoft \cite{microsoft2017ms17010}) reducen la susceptibilidad general de la red, pero son insuficientes ante vulnerabilidades de día cero o equipos huérfanos de mantenimiento.

Por ende, las estrategias reactivas automáticas de contención resultan indispensables. Lahari y Kumar \cite{lahari2025optimal} analizan estrategias óptimas de defensa frente a ransomware utilizando enfoques reactivos de parcheo y aislamiento. Concluyen que la inmunización selectiva e inmediata de los nodos sanos circundantes a un brote detiene la propagación exponencial con un coste de inoperabilidad significativamente menor al que causaría un colapso generalizado. Esto fundamenta el uso de protocolos automatizados de contingencia basados en el aislamiento lógico de los subsegmentos de red y la aplicación en tiempo de ejecución de restricciones de puertos y parches lógicos de emergencia, impidiendo que el ransomware salte entre zonas de la intranet universitaria.
